% Code excerpts follow the repository implementation; see docs/论文维护说明.md.
\clearpage
\appendix
\renewcommand{\theHsection}{appendix.\Alph{section}}
\renewcommand{\thesubsection}{\Alph{section}.\arabic{subsection}}
\renewcommand{\thetable}{\Alph{section}.\arabic{table}}
\definecolor{codeKeyword}{HTML}{225EA8}
\definecolor{codeString}{HTML}{98613D}
\definecolor{codeComment}{HTML}{557368}
\lstset{language=Python,basicstyle=\small\ttfamily\linespread{1.05}\selectfont,
 keywordstyle=\color{codeKeyword},commentstyle=\color{codeComment},stringstyle=\color{codeString},
 frame=single,rulecolor=\color{black!30},breaklines=true,breakatwhitespace=true,
 columns=fullflexible,keepspaces=true,showstringspaces=false,tabsize=4,
 aboveskip=8pt,belowskip=8pt}
\BeforeBeginEnvironment{lstlisting}{\Needspace{5\baselineskip}}
\section{补充结果与复现材料}
\label{sec:appendix-reproduction}

\subsection{问题一：特征维度与典型样本对应}
\subsubsection*{（1）声学与面部特征的维度定义}
文本表示为冻结BERT的768维末层向量。声学特征的列序、计算含义和可用条件见表~\ref{tab:appendix-audio-dimensions}；
13个MFCC与四个声学量共同组成17维表示。表中列序从1开始，数组索引从0开始。
静音窗仍保留声学特征，基频不可测时仅将该维的观测掩码置为假；
词级基频汇聚仅使用可测窗，并重新归一化交叠权重。
\begin{table}[htbp]
\centering
\caption{17维声学特征的排列与有效性}\label{tab:appendix-audio-dimensions}
\begin{tblr}{width=\linewidth,colspec={c c c X[c] X[c]},cells={valign=m,font=\normalsize\songti},
 colsep=4pt,rowsep=3pt,hline{1,7}={0.8pt},hline{2}={0.4pt}}
列序号 & 特征 & 维数 & 定义 & 可用条件 \\
1--13 & MFCC $c_0$--$c_{12}$ & 13 & 26通道对数mel谱的前13个正交倒谱系数 & 存在完整声学窗 \\
14 & 对数RMS & 1 & 去均值波形均方根幅值的自然对数 & 存在完整声学窗 \\
15 & 基频/Hz & 1 & 60--400 Hz内自相关峰值对应频率 & 峰值至少0.30，RMS至少$10^{-3}$ \\
16 & 频谱质心/Hz & 1 & 幅值谱的频率加权均值 & 完整声学窗；全零谱取零 \\
17 & 过零率 & 1 & 相邻采样点符号变化次数除以399 & 存在完整声学窗 \\
\end{tblr}
\end{table}


面部系数按英文名称排序，分组与列序见表~\ref{tab:appendix-video-dimensions}。
这些系数描述面部形变，其中基准系数\texttt{\_neutral}不代表情感分类中的中性概率。
词级面部特征仅汇聚单脸有效帧。电子结果中的特征维度表列出全部17个声学量与52个面部系数的名称、索引和掩码，
与每份特征文件内的名称序列一致。
\begin{table}[htbp]
\centering
\caption{52维面部系数的分组与排列}\label{tab:appendix-video-dimensions}
\begin{tabular*}{\linewidth}{@{\extracolsep{\fill}}ccc@{}}
\toprule
列序号 & 系数组 & 维数 \\
\midrule
1 & 基准系数 & 1 \\
2--6 & 眉部 & 5 \\
7--9 & 面颊 & 3 \\
10--23 & 眼部 & 14 \\
24--27 & 下颌 & 4 \\
28--50 & 口部 & 23 \\
51--52 & 鼻部 & 2 \\
\bottomrule
\end{tabular*}
\end{table}

\FloatBarrier

\subsubsection*{（2）典型样本逐词对应}
样本片段\texttt{-egA8-b7-3M\_9}的10个词均通过词时间筛查。
表~\ref{tab:appendix-word-times} 给出自动时间与交叠来源数量；试听确认内容及大致位置，
没有单独标注精确起止边界。
\begin{table}[htbp]
\centering
\caption{典型样本词时间与来源数量}\label{tab:appendix-word-times}
\begin{tabular*}{\linewidth}{@{\extracolsep{\fill}}ccccc@{}}
\toprule
序号 & 单词 & 时间区间/s & 语音窗数 & 视频帧数\\\midrule
1 & Secondly, & $[0.560, 1.040)$ & 51 & 16 \\
2 & another & $[1.260, 1.620)$ & 38 & 12 \\
3 & goal & $[1.680, 1.880)$ & 23 & 7 \\
4 & of & $[1.960, 2.000)$ & 6 & 2 \\
5 & exhibiting & $[2.100, 2.580)$ & 51 & 15 \\
6 & is & $[2.740, 2.800)$ & 9 & 2 \\
7 & building & $[2.900, 3.280)$ & 41 & 12 \\
8 & your & $[3.320, 3.460)$ & 16 & 5 \\
9 & company & $[3.520, 3.900)$ & 40 & 12 \\
10 & image. & $[4.040, 4.300)$ & 28 & 8 \\
\bottomrule
\end{tabular*}
\end{table}

\subsection{问题二：模型组件与逐类表现}
\subsubsection*{（1）候选模型的组件配置}
表~\ref{tab:appendix-model-components}列出统一训练器中六类候选的结构与辅助目标，
与正文表~\ref{tab:p2-candidate-comparison}的对应模型名称一致。

含文本的候选均采用冻结BERT，音视频卷积与BiLSTM候选仅使用音视频输入；
各模型均有极性分类与强度回归目标，回归项权重为0.2。
表中分支监督系数作用于相应分支的分类与回归联合损失。
跨模态注意力使用两层、四头、128维编码；MoE的两个专家输出通过样本级门控融合。
种子1、2、3依次为2026、2027、2028，各模型的训练轮次独立按验证集综合分数选择。
这些候选涉及编码结构、融合方式或辅助目标的变化，比较范围以表中配置为准。

\begin{table}[htbp]
\centering
\caption{候选模型的结构与辅助目标}\label{tab:appendix-model-components}
\begin{tblr}{width=\linewidth,colspec={X[1.4,c] X[1.1,c] X[1.2,c] X[1.2,c] X[1.2,c] c},
 cells={valign=m,font=\normalsize\songti},colsep=3pt,rowsep=2pt,hline{1,8}={0.8pt},hline{2}={0.4pt}}
模型 & 音视频编码 & 融合方式 & 分支监督 & 辅助约束 & 种子 \\
音视频卷积 & \makecell[c]{64维\\残差卷积} & 拼接与MLP & 无 & 无 & 1 \\
\makecell[c]{音视频\\BiLSTM} & 双向LSTM & 拼接与MLP & 无 & 无 & 1 \\
\makecell[c]{掩码跨模态\\注意力} & \makecell[c]{64维\\残差卷积} & \makecell[c]{两层跨模态\\注意力} & 无 & 无 & 1 \\
基础MoE & \makecell[c]{64维\\残差卷积} & 双专家门控 & \makecell[c]{文本0.2\\音视频0.3} & 无 & 1、2、3 \\
中性辅助MoE & \makecell[c]{64维\\残差卷积} & 双专家门控 & \makecell[c]{文本0.2\\音视频0.3} & \makecell[c]{中性二元\\损失0.3} & 1、2 \\
\makecell[c]{文本表征引导\\MoE} & \makecell[c]{64维\\残差卷积} & 双专家门控 & \makecell[c]{文本0.2\\音视频0.3} & \makecell[c]{文本重构\\0.02} & 1、2、3 \\
\end{tblr}
\end{table}

\FloatBarrier

\subsubsection*{（2）完整输入与缺失场景的逐类指标}
正文采用ACC、Macro-F1和回归指标比较缺失场景；表~\ref{tab:appendix-class-f1}
补充完整输入及21种缺失场景的逐类F1与加权F1。所有结果来自同一固定集成模型和728条验证样本。
\begingroup
\renewcommand{\arraystretch}{1.25}
\begin{longtable}{@{\extracolsep{\fill}}ccccccc@{}}
\caption{完整输入与缺失场景的逐类F1}\label{tab:appendix-class-f1}\\
\toprule
缺失模态 & 比例 & 位置 & 负向F1 & 中性F1 & 正向F1 & 加权F1 \\
\midrule
\endfirsthead
\multicolumn{7}{c}{表\thetable\ （续）}\\
\toprule
缺失模态 & 比例 & 位置 & 负向F1 & 中性F1 & 正向F1 & 加权F1 \\
\midrule
\endhead
\bottomrule
\endfoot
完整 & -- & -- & $0.6856$ & $0.4778$ & $0.7162$ & $0.6473$ \\
文本 & 20\% & 首部 & $0.6579$ & $0.4543$ & $0.7011$ & $0.6265$ \\
文本 & 20\% & 中部 & $0.6223$ & $0.4590$ & $0.6876$ & $0.6114$ \\
文本 & 20\% & 尾部 & $0.6049$ & $0.4011$ & $0.6522$ & $0.5753$ \\
文本 & 40\% & 首部 & $0.6208$ & $0.4474$ & $0.6688$ & $0.5993$ \\
文本 & 40\% & 中部 & $0.6047$ & $0.4290$ & $0.6687$ & $0.5900$ \\
文本 & 40\% & 尾部 & $0.5996$ & $0.4113$ & $0.6483$ & $0.5746$ \\
文本 & 100\% & 全部 & $0.3850$ & $0.3123$ & $0.5340$ & $0.4358$ \\
语音 & 20\% & 首部 & $0.6840$ & $0.4765$ & $0.7124$ & $0.6447$ \\
语音 & 20\% & 中部 & $0.6872$ & $0.4820$ & $0.7162$ & $0.6488$ \\
语音 & 20\% & 尾部 & $0.6825$ & $0.4791$ & $0.7170$ & $0.6471$ \\
语音 & 40\% & 首部 & $0.6840$ & $0.4765$ & $0.7124$ & $0.6447$ \\
语音 & 40\% & 中部 & $0.6841$ & $0.4780$ & $0.7124$ & $0.6451$ \\
语音 & 40\% & 尾部 & $0.6825$ & $0.4847$ & $0.7170$ & $0.6485$ \\
语音 & 100\% & 全部 & $0.6975$ & $0.4903$ & $0.7108$ & $0.6513$ \\
视频 & 20\% & 首部 & $0.6825$ & $0.4807$ & $0.7173$ & $0.6476$ \\
视频 & 20\% & 中部 & $0.6857$ & $0.4807$ & $0.7122$ & $0.6462$ \\
视频 & 20\% & 尾部 & $0.6824$ & $0.4848$ & $0.7126$ & $0.6465$ \\
视频 & 40\% & 首部 & $0.6857$ & $0.4848$ & $0.7162$ & $0.6491$ \\
视频 & 40\% & 中部 & $0.6857$ & $0.4807$ & $0.7122$ & $0.6462$ \\
视频 & 40\% & 尾部 & $0.6792$ & $0.4877$ & $0.7106$ & $0.6454$ \\
视频 & 100\% & 全部 & $0.6634$ & $0.4736$ & $0.6962$ & $0.6306$ \\
\end{longtable}
\endgroup

\subsection{问题三：模态贡献与局部证据}
\subsubsection*{（1）极性概率贡献}
表~\ref{tab:appendix-contributions} 补充附件四全部20条样本的有符号Shapley贡献。
贡献针对各样本完整输入的预测类别计算，正值表示提高该类别概率，负值表示降低；
三模态之和等于完整输入与全遮挡输入的概率差。数值保留四位小数。
\begingroup
\renewcommand{\arraystretch}{1.25}
\begin{longtable}{@{\extracolsep{\fill}}ccccc@{}}
\caption{附件四全部样本的三模态贡献}\label{tab:appendix-contributions}\\
\toprule
编号 & 文本贡献 & 语音贡献 & 视频贡献 & 极性主模态 \\
\midrule
\endfirsthead
\multicolumn{5}{c}{表\thetable\ （续）}\\
\toprule
编号 & 文本贡献 & 语音贡献 & 视频贡献 & 极性主模态 \\
\midrule
\endhead
\bottomrule
\endfoot
01 & $0.1652$ & $0.0421$ & $-0.1003$ & 文本 \\
02 & $0.1711$ & $-0.0794$ & $0.0561$ & 文本 \\
03 & $0.2647$ & $-0.1015$ & $0.0421$ & 文本 \\
04 & $0.5303$ & $-0.0892$ & $-0.0444$ & 文本 \\
05 & $0.1422$ & $0.0152$ & $0.1892$ & 视频 \\
06 & $0.5870$ & $-0.0262$ & $-0.0482$ & 文本 \\
07 & $0.3776$ & $0.0481$ & $0.1177$ & 文本 \\
08 & $0.5085$ & $-0.0670$ & $-0.0093$ & 文本 \\
09 & $0.5833$ & $0.0362$ & $0.0829$ & 文本 \\
10 & $0.6612$ & $-0.0445$ & $0.0190$ & 文本 \\
11 & $0.2944$ & $-0.0623$ & $0.1264$ & 文本 \\
12 & $0.5900$ & $-0.1390$ & $0.0308$ & 文本 \\
13 & $0.1220$ & $0.0481$ & $0.0000$ & 文本 \\
14 & $0.4113$ & $0.0458$ & $-0.2255$ & 文本 \\
15 & $0.4889$ & $0.1144$ & $0.0771$ & 文本 \\
16 & $0.7318$ & $-0.0585$ & $0.0241$ & 文本 \\
17 & $0.5789$ & $0.0210$ & $0.1000$ & 文本 \\
18 & $0.1601$ & $-0.0499$ & $-0.0108$ & 文本 \\
19 & $0.2497$ & $-0.1123$ & $0.0739$ & 文本 \\
20 & $0.4820$ & $0.0219$ & $0.1167$ & 文本 \\
\end{longtable}
\endgroup


\subsubsection*{（2）情感强度贡献}
表~\ref{tab:appendix-strength-contributions}给出符号投影前原始回归输出的Shapley贡献。
正值表示该模态提高原始强度，负值表示降低原始强度；强度主模态按贡献绝对值最大者确定。
其方向与数值均针对连续回归输出。预测为负向时，负贡献也可能增强模型的负向判断；
贡献绝对值描述影响幅度。中性样本的最终强度经投影置零，其原始回归贡献仍可非零。
% Generated by src/q3_explanation/export_paper_tables.py
\begingroup
\setlength{\tabcolsep}{3pt}
\begin{longtable}{@{\extracolsep{\fill}}ccccc@{}}
\caption{附件四全部样本的原始强度贡献}\label{tab:appendix-strength-contributions}\\
\toprule
编号 & 文本贡献 & 音频贡献 & 视频贡献 & 强度主模态 \\
\midrule
\endfirsthead
\multicolumn{5}{c}{表 \thetable\ （续）}\\
\toprule
编号 & 文本贡献 & 音频贡献 & 视频贡献 & 强度主模态 \\
\midrule
\endhead
\bottomrule
\endfoot
01 & $0.2221$ & $0.2248$ & $0.0123$ & 音频 \\
02 & $0.2777$ & $-0.0892$ & $0.1061$ & 文本 \\
03 & $-0.4741$ & $0.2458$ & $0.0935$ & 文本 \\
04 & $-0.9951$ & $0.3619$ & $0.4530$ & 文本 \\
05 & $0.3565$ & $0.1142$ & $0.2934$ & 文本 \\
06 & $0.9602$ & $0.1449$ & $-0.2097$ & 文本 \\
07 & $0.8022$ & $0.1899$ & $0.1775$ & 文本 \\
08 & $1.0130$ & $-0.0382$ & $-0.1147$ & 文本 \\
09 & $-1.7384$ & $0.0636$ & $-0.0189$ & 文本 \\
10 & $-1.6117$ & $0.1736$ & $0.1695$ & 文本 \\
11 & $-0.2056$ & $0.0320$ & $-0.0754$ & 文本 \\
12 & $-0.9900$ & $0.2634$ & $0.0023$ & 文本 \\
13 & $0.5563$ & $-0.1280$ & $0.0000$ & 文本 \\
14 & $-0.1583$ & $0.2496$ & $0.4426$ & 视频 \\
15 & $1.0617$ & $0.1762$ & $0.1476$ & 文本 \\
16 & $-1.8590$ & $0.1063$ & $0.0878$ & 文本 \\
17 & $1.3118$ & $0.1801$ & $0.0404$ & 文本 \\
18 & $-0.3351$ & $0.3012$ & $-0.0153$ & 文本 \\
19 & $-0.1049$ & $0.2025$ & $0.1093$ & 音频 \\
20 & $0.9210$ & $-0.0027$ & $0.2418$ & 文本 \\
\end{longtable}
\endgroup


\subsubsection*{（3）各模态的关键局部窗口}
表~\ref{tab:appendix-local-evidence} 补充三种模态各自概率下降最大的局部窗口，位置采用
特征序列零基索引，区间包含两端。下降量为完整输入概率减去遮挡后概率，负值表示遮挡后概率上升。
表中区间表示特征位置，原素材的时间对应见正文表~\ref{tab:p3-attachment4-evidence}。
\begingroup
\renewcommand{\arraystretch}{1.25}
\begin{longtable}{@{\extracolsep{\fill}}ccccccc@{}}
\caption{附件四各模态的高分局部窗口}\label{tab:appendix-local-evidence}\\
\toprule
编号 & 文本位置 & 下降量 & 语音位置 & 下降量 & 视频位置 & 下降量 \\
\midrule
\endfirsthead
\multicolumn{7}{c}{表\thetable\ （续）}\\
\toprule
编号 & 文本位置 & 下降量 & 语音位置 & 下降量 & 视频位置 & 下降量 \\
\midrule
\endhead
\bottomrule
\endfoot
01 & $[7, 10]$ & $0.1329$ & $[11, 14]$ & $0.0148$ & $[1, 4]$ & $0.0051$ \\
02 & $[9, 12]$ & $0.2183$ & $[11, 14]$ & $0.0049$ & $[2, 5]$ & $0.0016$ \\
03 & $[17, 23]$ & $0.2603$ & $[22, 28]$ & $0.0116$ & $[24, 30]$ & $0.0064$ \\
04 & $[21, 25]$ & $0.3387$ & $[3, 7]$ & $0.0086$ & $[10, 14]$ & $0.0032$ \\
05 & $[1, 3]$ & $0.2309$ & $[3, 5]$ & $0.0121$ & $[2, 4]$ & $0.0199$ \\
06 & $[1, 7]$ & $0.4279$ & $[21, 27]$ & $0.0081$ & $[24, 30]$ & $0.0038$ \\
07 & $[21, 30]$ & $0.3356$ & $[16, 25]$ & $0.0050$ & $[38, 47]$ & $0.0049$ \\
08 & $[2, 3]$ & $0.2942$ & $[8, 9]$ & $0.0031$ & $[10, 11]$ & $0.0049$ \\
09 & $[7, 10]$ & $0.0117$ & $[14, 17]$ & $0.0003$ & $[2, 5]$ & $0.0005$ \\
10 & $[12, 17]$ & $0.2466$ & $[13, 18]$ & $0.0014$ & $[3, 8]$ & $0.0012$ \\
11 & $[2, 4]$ & $0.4178$ & $[6, 8]$ & $0.0027$ & $[2, 4]$ & $0.0031$ \\
12 & $[34, 41]$ & $0.1515$ & $[33, 40]$ & $0.0094$ & $[1, 8]$ & $0.0037$ \\
13 & $[23, 28]$ & $0.0718$ & $[9, 14]$ & $0.0165$ & -- & -- \\
14 & $[21, 27]$ & $0.1428$ & $[24, 30]$ & $0.0118$ & $[1, 7]$ & $0.0006$ \\
15 & $[5, 8]$ & $0.3485$ & $[1, 4]$ & $0.0008$ & $[18, 21]$ & $0.0019$ \\
16 & $[9, 10]$ & $0.1679$ & $[9, 10]$ & $0.0009$ & $[5, 6]$ & $0.0004$ \\
17 & $[10, 14]$ & $0.0713$ & $[18, 22]$ & $0.0018$ & $[4, 8]$ & $0.0020$ \\
18 & $[3, 12]$ & $0.2233$ & $[1, 10]$ & $0.0079$ & $[1, 10]$ & $0.0025$ \\
19 & $[12, 19]$ & $0.2364$ & $[14, 21]$ & $0.0082$ & $[30, 37]$ & $0.0084$ \\
20 & $[19, 27]$ & $0.1387$ & $[3, 11]$ & $0.0019$ & $[13, 21]$ & $0.0074$ \\
\end{longtable}
\endgroup

\subsection{工具、预训练模型与实验配置}
\label{sec:appendix-runtime}
问题一按样本片段编号保存100份压缩NPZ文件。
每份文件包含官方文本与标签、原生序列、词级特征、实际长度、时间与观测掩码，以及来源索引和聚合权重；
另提供按样本和模态汇总的300行结果表、异常统计与逐词质量记录。
复现材料保留模型修订号、权重校验值、依赖版本、解码参数及听审原始记录。

表~\ref{tab:appendix-tools} 统一列出建模所用工具、公开预训练模型及核心设置。
公开权重包括google-bert/bert-base-uncased、facebook/wav2vec2-base-960h和Face Landmarker模型包，
分别用于文本编码、语音时间对齐和面部特征提取，使用过程中不更新这些权重。
混合专家模型中的门控网络、预测头和重构头使用附件二训练集学习参数；验证集用于选择模型结构与训练轮次。
表~\ref{tab:appendix-entry}列出核心程序及其输入，详细运行步骤见电子附件中的复现说明。
环境由\texttt{pixi.toml}和\texttt{pixi.lock}固定；训练使用种子2026、2027、2028，推理采用已选定的模型参数。
BERT与wav2vec 2.0使用权重修订号标识版本；Face Landmarker使用官方模型包版本和文件SHA-256标识。

\noindent\textbf{BERT 完整权重修订号：}\par\nopagebreak
\noindent\texttt{86b5e0934494bd15c9632b12f734a8a67f723594}\par\smallskip
\noindent\textbf{wav2vec 2.0 完整权重修订号：}\par\nopagebreak
\noindent\texttt{22aad52d435eb6dbaf354bdad9b0da84ce7d6156}\par\smallskip
\noindent\textbf{Face Landmarker 模型包：}\texttt{float16/1}\par\nopagebreak
\noindent\textbf{完整 SHA-256：}\par\nopagebreak
\noindent{\small\texttt{64184e229b263107bc2b804c6625db1341ff2bb731874b0bcc2fe6544e0bc9ff}}\par

\begingroup
\normalsize\songti
\setlength{\tabcolsep}{3pt}
\begin{longtable}{@{}C{36mm}C{35mm}C{55mm}C{\dimexpr\linewidth-126mm-6\tabcolsep\relax}@{}}
\caption{开源工具与预训练模型的版本、参数及用途}\label{tab:appendix-tools}\\
\toprule
工具或模型 & 版本或权重标识 & 核心设置 & 使用场景 \\
\midrule
\endfirsthead
\multicolumn{4}{c}{表\thetable\ （续）}\\
\toprule
工具或模型 & 版本或权重标识 & 核心设置 & 使用场景 \\
\midrule
\endhead
\bottomrule
\endfoot
\makecell[c]{BERT\\bert-base-uncased} & \makecell[c]{Transformers 4.57.6\\修订号86b5e093} & 冻结权重；词元表示768维；预测模型拼接CLS与观测词元均值得1536维表示 & 三问文本编码 \\
\makecell[c]{wav2vec 2.0\\wav2vec2-base-960h} & \makecell[c]{Transformers 4.57.6\\修订号22aad52d} & 冻结英语CTC模型；16 kHz音频；输出步长20 ms；字符置信度阈值0.50 & 问题一词对齐、问题三时间候选 \\
\makecell[c]{MediaPipe\\Face Landmarker} & \makecell[c]{MediaPipe 1.0.1\\模型包float16/1} & VIDEO模式；帧宽640像素；最多检测2张脸；检测、存在与跟踪阈值均0.5；输出52维系数 & 问题一视频特征 \\
PyTorch & 2.11.0+cu128 & AdamW；学习率$3\times10^{-4}$；batch=64；随机种子2026、2027、2028；其余配置见表~\ref{tab:p2-training-config} & 情感模型训练、推理与遮挡解释 \\
FFmpeg / ffprobe & 4.4.2 & 保留解码时间基准；音频转16 kHz单声道；读取帧显示时间与时长 & 原始音视频解码与时间记录 \\
NumPy / SciPy & 2.5.3 / 1.18.1 & 声学窗25 ms、步长10 ms；13维MFCC及能量、过零率、质心、基频，共17维 & 问题一声学特征与区间聚合 \\
scikit-learn & 1.9.1 & 固定三类标签；计算ACC、Macro-F1、逐类F1及MAE & 验证集指标评估 \\
Matplotlib & 3.11.2 & 宋体与Times New Roman；矢量PDF；统一模态颜色与指标精度 & 论文结果可视化 \\
Python & 3.12.14 & 使用锁定依赖环境；预训练工具按推理模式运行 & 三问程序运行 \\
\end{longtable}
\endgroup

\begin{table}[!htbp]
\centering
\caption{核心程序及其输入}\label{tab:appendix-entry}
\begin{tabularx}{\linewidth}{@{}C{23mm}C{65mm}Y@{}}
\toprule
计算任务 & 代码位置 & 输入数据与模型 \\
\midrule
特征提取 & \path{src/q1_features/} & 附件一视频、转写及预训练提取器 \\
训练与评估 & \path{src/q2_prediction/} & 附件二对齐版训练集与验证集 \\
专项预测 & \path{src/q2_prediction/} & 附件三对齐版、三种子参数 \\
预测解释 & \path{src/q3_explanation/} & 附件四对齐版及视频、三种子参数 \\
\bottomrule
\end{tabularx}
\end{table}

\FloatBarrier
\subsection{结果字段说明}
正文已给出100条特征汇总、30条专项预测与20条专项解释。
表~\ref{tab:appendix-fields} 补充电子结果中的字段含义，便于读取完整结果。
\begin{table}[!htbp]
\centering
\caption{电子结果的主要字段说明}\label{tab:appendix-fields}
\begin{tabularx}{\linewidth}{@{}C{25mm}C{36mm}Y@{}}
\toprule
结果材料 & 字段组 & 内容与含义 \\
\midrule
特征NPZ & 原生特征与位置 & 文本词元、声学短窗、面部系数及各自的来源位置 \\
特征NPZ & 词级特征与掩码 & 词起止时间、结构有效掩码、模态来源与面部观测掩码 \\
特征NPZ & 来源索引与权重 & 从原生特征重算词级值所需的索引和聚合权重 \\
附件三CSV & 编号、极性、强度 & 全部30条样本的预测，强度按类别符号规则投影 \\
附件四解释 & 模态贡献 & 三模态有符号贡献、正贡献占比及最大贡献模态 \\
附件四解释 & 局部窗口 & 各模态窗口范围与预测类别概率下降量 \\
内容核查 & 时间与对应关系 & 自动时间估计、音文一致性及视频画面对应的核查结论 \\
\bottomrule
\end{tabularx}
\end{table}

\FloatBarrier
\Needspace{8\baselineskip}
\section{核心代码}
\label{sec:appendix-code}
本附录列出词级对齐、三种子MoE集成和八组合解释的核心函数，完整实现见电子附件。


\subsection{特征提取与时序对齐}
\subsubsection*{（1）文本编码与词元对应}
以下函数采用冻结BERT提取词元表示，并保存原文词与子词的对应索引。
该索引用于关联文本内容；词的起止时间由另外的对齐过程确定。

% Source: src/q1_features/feature_extractors.py :: encode_text
\begin{lstlisting}
def encode_text(text: str, tokenizer, model, device) -> dict:
    import torch

    # Pretokenize by official whitespace words so word_ids refer to the same
    # word indices as the CTC alignment and human review tables.
    words = text.split()
    encoding = tokenizer(words, is_split_into_words=True, add_special_tokens=True, truncation=True,
                         max_length=MAX_TEXT_TOKENS, return_tensors="pt")
    input_ids = encoding["input_ids"][0]
    attention = encoding["attention_mask"][0]
    token_type = encoding.get("token_type_ids", torch.zeros_like(input_ids))[0]
    n_before = len(tokenizer(words, is_split_into_words=True, add_special_tokens=True,
                             truncation=False)["input_ids"])
    with torch.no_grad():
        outputs = model(input_ids=input_ids.unsqueeze(0).to(device),
                        attention_mask=attention.unsqueeze(0).to(device))
    hidden = outputs.last_hidden_state[0].cpu().numpy().astype(np.float32)
    # 词 -> 子词映射（用快速分词器的 word_ids）
    word_ids = encoding.word_ids(0) if hasattr(encoding, "word_ids") else None
    if word_ids is None:
        word_ids = tokenizer(words, is_split_into_words=True, add_special_tokens=True,
                             truncation=True, max_length=MAX_TEXT_TOKENS).word_ids()
    mapping = [-1 if w is None else int(w) for w in word_ids]
    return {
        "input_ids": input_ids.numpy().astype(np.int32),
        "attention_mask": attention.numpy().astype(bool),
        "token_type_ids": token_type.numpy().astype(np.int32),
        "last_hidden_state": hidden,
        "cls_vector": hidden[0].copy(),
        "word_ids": np.asarray(mapping, dtype=np.int32),
        "n_tokens": int(input_ids.shape[0]),
        "n_tokens_before_truncation": int(n_before),
        "truncated": bool(n_before > MAX_TEXT_TOKENS),
    }
\end{lstlisting}

\subsubsection*{（2）CTC最大概率路径}
以下函数在给定目标字符序列下求解CTC路径，并返回每个字符的首帧、末帧后一位置及置信度。
\texttt{lp}为声学模型的对数后验，空白符编号为0；相邻相同字符必须经过空白状态。
字符置信度取其占用帧后验的几何均值。
% Generated from the actual Q1 source; formatting only.
\begin{lstlisting}
def min_required_frames(targets: list[int]) -> int:
    """CTC 强制对齐所需最少帧数：每个目标字符至少 1 帧，且相邻相同字符之间必须有 blank。"""
    repeats = sum((1 for i in range(1, len(targets)) if targets[i] == targets[i - 1]))
    return len(targets) + repeats

def ctc_path(lp, targets):
    """Vectorized Viterbi; return first AND last emitting frames, excluding blanks."""
    T, V = lp.shape
    if not targets or T < min_required_frames(targets):
        raise ValueError('Unreachable CTC target')
    labels = np.zeros(2 * len(targets) + 1, dtype=np.int64)
    labels[1::2] = targets
    E = len(labels)
    neg = -np.inf
    previous = np.full(E, neg)
    previous[:2] = lp[0, labels[:2]]
    back = np.zeros((T, E), np.int8)
    skip = np.zeros(E, bool)
    skip[2:] = (labels[2:] != 0) & (labels[2:] != labels[:-2])
    for t in range(1, T):
        advance = np.r_[neg, previous[:-1]]
        jump = np.r_[neg, neg, previous[:-2]]
        jump[~skip] = neg
        options = np.stack([previous, advance, jump])
        move = options.argmax(axis=0)
        back[t] = move
        previous = options[move, np.arange(E)] + lp[t, labels]
    state = E - 1 if previous[-1] >= previous[-2] else E - 2
    score = previous[state]
    if not np.isfinite(score):
        raise ValueError('No finite CTC path')
    path = np.empty(T, np.int32)
    for t in range(T - 1, -1, -1):
        path[t] = state
        state -= int(back[t, state])
    starts = []
    ends = []
    conf = []
    for j, token in enumerate(targets):
        frames = np.flatnonzero(path == 2 * j + 1)
        if not len(frames):
            raise ValueError('Missing emitting state')
        starts.append(int(frames[0]))
        ends.append(int(frames[-1] + 1))
        conf.append(float(np.exp(lp[frames, token].mean())))
    return (np.array(starts), np.array(ends), np.array(conf), float(score))
\end{lstlisting}


\subsubsection*{（3）字符路径到官方词边界}
规范化文本与官方词保持对应。下列函数将目标字符位置映射到各词，
以20 ms输出步长换算起止时间，并对词内非分隔字符的置信度取平均。
\texttt{offset}为局部音频在完整音频中的输出帧偏移。
% Generated from the actual Q1 source; formatting only.
\begin{lstlisting}
STEP = 0.02

def targets_for(norms, vocab):
    tokens = []
    positions = []
    for word in norms:
        if tokens:
            tokens.append(vocab['|'])
        indexes = []
        for c in word:
            token = vocab['|'] if c == ' ' else vocab[c]
            indexes.append(len(tokens))
            tokens.append(token)
        positions.append(indexes)
    return (tokens, positions)

def align(lp, norms, vocab, offset=0):
    targets, positions = targets_for(norms, vocab)
    starts, ends, confidence, score = ctc_path(lp, targets)
    rows = []
    for pos in positions:
        letters = [j for j in pos if targets[j] != vocab['|']]
        if not letters:
            rows.append(None)
            continue
        rows.append(dict(start=round((offset + starts[letters[0]]) * STEP, 6), end=round((offset + ends[letters[-1]]) * STEP, 6), confidence=float(confidence[letters].mean()), confidence_min=float(confidence[letters].min())))
    return (rows, score)
\end{lstlisting}


\subsubsection*{（4）局部边界复核与词时间筛选}
\texttt{blocks}给出内容核查确定的对齐范围，\texttt{norms}为规范化词序列，
\texttt{results}保留全部官方词位置，\texttt{silence}为检测得到的静音区间。
在每词前后各三个词、时间两侧各160 ms的局部范围内重新求解路径，
结合字符置信度、边界变化、静音、跨度及交叠情况确定可用词时间。
% Generated from the actual Q1 source; formatting only.
\begin{lstlisting}
CONFIDENCE = 0.50
STABILITY = 0.10

for block in blocks:
    ids = block['indices']
    left, right = (block['left'], block['right'])
    ns = [norms[i] for i in ids]
    try:
        global_rows, _ = align(lp[left:right], ns, vocab, left)
    except ValueError:
        for i in ids:
            results[i]['reason'] = 'ctc_path_unreachable'
        continue
    for j, i in enumerate(ids):
        rr = results[i]
        g = global_rows[j]
        if not g:
            rr['reason'] = 'no_supported_target_characters'
            continue
        rr.update(candidate_start=g['start'], candidate_end=g['end'], confidence=g['confidence'], confidence_min=g['confidence_min'])
        a = max(0, j - 3)
        b = min(len(ids), j + 4)
        nonempty = [v for v in global_rows[a:b] if v]
        lo = max(left, round(nonempty[0]['start'] / STEP) - 8)
        hi = min(right, round(nonempty[-1]['end'] / STEP) + 8)
        try:
            local, _ = align(lp[lo:hi], [norms[x] for x in ids[a:b]], vocab, lo)
            v = local[j - a]
            if not v:
                raise ValueError('No local word')
        except ValueError:
            rr['reason'] = 'local_path_unreachable'
            continue
        shift = max(abs(g['start'] - v['start']), abs(g['end'] - v['end']))
        rr.update(local_start=v['start'], local_end=v['end'], local_confidence=v['confidence'], boundary_shift_s=round(shift, 6))
        entirely_silent = any((v['start'] >= z['start'] and v['end'] <= (z['end'] if z['end'] is not None else len(lp) * STEP) for z in silence))
        duration = v['end'] - v['start']
        max_duration = max(1.0, 0.2 * len(norms[i].replace(' ', '')))
        if min(g['confidence'], v['confidence']) < CONFIDENCE:
            reason = 'low_character_confidence'
        elif shift > STABILITY + 1e-08:
            reason = 'unstable_local_boundary'
        elif entirely_silent:
            reason = 'inside_detected_silence'
        elif duration > max_duration:
            reason = 'excessive_word_span'
        elif not 0 <= v['start'] < v['end'] <= len(lp) * STEP + 1e-08:
            reason = 'invalid_time'
        else:
            reason = 'accepted_automatic_with_content_evidence'
            rr.update(accepted=True, start=v['start'], end=v['end'])
        rr['reason'] = reason

previous = None

for rr in results:
    if not rr['accepted']:
        continue
    if previous and rr['start'] < previous['end'] - 1e-08:
        for x in [previous, rr]:
            x.update(accepted=False, start=None, end=None, reason='overlap_after_local_alignment')
        previous = None
    else:
        previous = rr
\end{lstlisting}


\subsubsection*{（5）词区间与原生语音窗聚合}
交叠长度决定语音窗权重，基频仅使用可测窗并重新归一化。
\begin{lstlisting}
PITCH_DIM = 14
def overlap_sources(start, end, left, right):
    overlap = np.maximum(0., np.minimum(right, end) - np.maximum(left, start))
    indices = np.flatnonzero(overlap > 1e-9)
    lengths = overlap[indices]
    return indices, lengths, lengths / lengths.sum() if len(lengths) else lengths


def pooled_audio(features, pitch_mask, indices, weights):
    out = np.zeros(features.shape[1], dtype=np.float64)
    if len(indices):
        out = weights @ features[indices].astype(np.float64)
        voiced = pitch_mask[indices]
        out[PITCH_DIM] = (weights[voiced] @ features[indices[voiced], PITCH_DIM]
                          / weights[voiced].sum()) if voiced.any() else 0.
    return out


\end{lstlisting}
\subsubsection*{（6）视频显示区间与词级汇聚}
视频帧的显示区间由真实时间戳和帧时长确定。多人脸帧保留原生记录，
仅单脸有效帧参与词级面部汇聚。
\begin{lstlisting}
def display_intervals(tb):
    v = tb['video_frames']
    start = np.array(v['display_times'], float)
    duration = np.array(v['pkt_durations'], float)
    stop = np.r_[start[1:], float(v['coverage_end_seconds'])]
    end = np.where(
        np.isfinite(duration) & (duration > 0),
        np.minimum(start + duration, stop), stop)
    assert np.all(np.diff(start) > 0) and np.all(end > start)
    return start, end

# For each accepted word [start[i], end[i]):
vi, overlap, _ = overlap_sources(
    start[i], end[i], pts, frame_end)
video_source_observed[i] = len(vi) > 0
weights = overlap * single_face_valid[vi]
if weights.sum() > 0:
    weights /= weights.sum()
    word_video[i] = weights @ frame_features[vi]
    face_observed[i] = True
# Save vi, overlap and weights for independent reconstruction.
\end{lstlisting}

\subsection{局部缺失预测模型}
\subsubsection*{（1）连续缺失区间}
以下公共函数在已观测位置上选取连续区间，供固定位置评估与随机训练遮挡调用。
区间长度按观测位置数和移除比例确定。

% Source: src/common/prep.py :: missing_span
\begin{lstlisting}
def missing_span(observed: torch.Tensor, fraction: float, position: str,
                 rng: random.Random | None = None) -> torch.Tensor:
    """Choose a contiguous interval in the valid observed prefix of each sample."""
    missing = torch.zeros_like(observed)
    for row in range(len(observed)):
        positions = torch.nonzero(observed[row], as_tuple=True)[0]
        n = len(positions)
        if n == 0:
            continue
        width = max(1, min(n, round(n * fraction)))
        if position == "start":
            start = 0
        elif position == "middle":
            start = (n - width) // 2
        elif position == "end":
            start = n - width
        elif position == "random" and rng is not None:
            start = rng.randint(0, n - width)
        else:
            raise ValueError(position)
        missing[row, positions[start:start + width]] = True
    return missing
\end{lstlisting}

\Needspace{7\baselineskip}
上述区间通过以下函数同时作用于输入与观测掩码，结构掩码保持不变。

% Source: src/common/prep.py :: apply_missing
\begin{lstlisting}
def apply_missing(batch: dict, modality: str, fraction: float, position: str,
                  rng: random.Random | None = None, rows: np.ndarray | None = None) -> None:
    key = f"{modality}_observed"
    missing = missing_span(batch[key], fraction, position, rng)
    if rows is not None:
        selected = torch.zeros(len(missing), device=missing.device, dtype=torch.bool)
        selected[rows] = True
        missing &= selected[:, None]
    batch[key] &= ~missing
    if modality == "text":
        batch["tokens"][missing] = 103  # [MASK], keeping the structural attention mask.
    else:
        batch[modality][missing] = 0
\end{lstlisting}

\subsubsection*{（2）音视频时序编码与混合专家模型}
音视频编码器先做特征投影与时序卷积，再组合掩码均值、注意力池化和观测比例。
本文模型使用64维编码表示。

% Source: src/common/encoders.py :: AVEncoder
\begin{lstlisting}
class AVEncoder(nn.Module):
    def __init__(self, input_dim: int, width: int = 96) -> None:
        super().__init__()
        self.project = nn.Sequential(nn.Linear(input_dim, width), nn.LayerNorm(width), nn.GELU())
        self.position = nn.Parameter(torch.zeros(1, 50, width))
        self.conv = nn.Conv1d(width, width, kernel_size=3, padding=1)
        self.score = nn.Linear(width, 1)

    def forward(self, x: torch.Tensor, observed: torch.Tensor) -> torch.Tensor:
        h = self.project(x) + self.position[:, :x.shape[1]]
        h = h * observed.unsqueeze(-1)
        h = torch.nn.functional.gelu(h + self.conv(h.transpose(1, 2)).transpose(1, 2))
        h = h * observed.unsqueeze(-1)
        count = observed.sum(1, keepdim=True).clamp_min(1)
        mean = h.sum(1) / count
        score = self.score(h).squeeze(-1).masked_fill(~observed, -1e4)
        active = observed.any(1)
        score[~active, 0] = 0
        attention = torch.softmax(score, dim=1)
        pooled = (attention.unsqueeze(-1) * h).sum(1)
        pooled = pooled * active.unsqueeze(-1)
        fraction = observed.float().mean(1, keepdim=True)
        return torch.cat([mean, pooled, fraction], dim=-1)
\end{lstlisting}

以下代码包含文本表征引导MoE及音视频单分支、BiLSTM和跨模态注意力对照。
主模型对应\texttt{gate\_recon}分支。

% Source: src/q2_prediction/experiment_q2_mask_fusion.py :: Specialist
\begin{lstlisting}
class Specialist(nn.Module):
    def __init__(self, mode: str):
        super().__init__()
        self.mode = mode
        encoder = MaskedBiLSTMEncoder if mode == "av_lstm" else lambda dim: AVEncoder(dim, 64)
        self.audio = encoder(74)
        self.vision = encoder(35)
        av_width = 129 * 2
        self.av_backbone = nn.Sequential(nn.Linear(av_width, 192), nn.LayerNorm(192),
                                         nn.GELU(), nn.Dropout(.25), nn.Linear(192, 128),
                                         nn.GELU())
        self.av_class = nn.Linear(128, 3)
        self.av_reg = nn.Linear(128, 1)
        if mode not in ("av", "av_lstm"):
            self.text_backbone = nn.Sequential(nn.Linear(1536, 192), nn.LayerNorm(192),
                                               nn.GELU(), nn.Dropout(.25),
                                               nn.Linear(192, 128), nn.GELU())
            self.text_class = nn.Linear(128, 3)
            self.text_reg = nn.Linear(128, 1)
            self.gate = nn.Sequential(nn.Linear(128 * 2 + 3, 64), nn.GELU(),
                                      nn.Linear(64, 1))
        if mode == "gate_recon":
            self.reconstruct = nn.Linear(128, 1536)
        if mode == "cross_attn":
            self.text_token = nn.Linear(1536, 128)
            self.audio_token = nn.Linear(129, 128)
            self.vision_token = nn.Linear(129, 128)
            layer = nn.TransformerEncoderLayer(128, 4, 256, .2, batch_first=True)
            self.cross_encoder = nn.TransformerEncoder(layer, 2)
            self.cross_class = nn.Linear(128, 3)
            self.cross_reg = nn.Linear(128, 1)

    def forward(self, batch: dict, bert: nn.Module | None):
        audio = self.audio(batch["audio"], batch["audio_observed"])
        vision = self.vision(batch["vision"], batch["vision_observed"])
        av_h = self.av_backbone(torch.cat([audio, vision], dim=-1))
        av_logits, av_reg = self.av_class(av_h), self.av_reg(av_h).squeeze(-1)
        if self.mode in ("av", "av_lstm"):
            return av_logits, av_reg, {"av_logits": av_logits, "av_reg": av_reg}
        with torch.no_grad():
            h = bert(input_ids=batch["tokens"], attention_mask=batch["struct"].long(),
                     token_type_ids=batch["segment"]).last_hidden_state.float()
            w = batch["text_observed"].unsqueeze(-1).float()
            text = torch.cat([h[:, 0], (h * w).sum(1) / w.sum(1).clamp_min(1)], dim=-1)
        if self.mode == "cross_attn":
            tokens = torch.stack([self.text_token(text), self.audio_token(audio),
                                  self.vision_token(vision)], dim=1)
            active = torch.stack([batch[f"{m}_observed"].any(1)
                                  for m in ("text", "audio", "vision")], dim=1)
            mixed = self.cross_encoder(tokens, src_key_padding_mask=~active)
            pooled = (mixed * active.unsqueeze(-1)).sum(1) / active.sum(1, keepdim=True).clamp_min(1)
            logits, reg = self.cross_class(pooled), self.cross_reg(pooled).squeeze(-1)
            return logits, reg, {"av_logits": av_logits, "av_reg": av_reg}
        text_h = self.text_backbone(text)
        text_logits, text_reg = self.text_class(text_h), self.text_reg(text_h).squeeze(-1)
        fractions = torch.stack([batch[f"{m}_observed"].float().mean(1)
                                 for m in ("text", "audio", "vision")], dim=-1)
        g = torch.sigmoid(self.gate(torch.cat([text_h, av_h, fractions], dim=-1))).squeeze(-1)
        av_available = (fractions[:, 1] + fractions[:, 2]) > 0
        text_available = fractions[:, 0] > 0
        g = torch.where(text_available, g, torch.zeros_like(g))
        g = torch.where(av_available, g, torch.ones_like(g))
        logits = g[:, None] * text_logits + (1 - g[:, None]) * av_logits
        reg = g * text_reg + (1 - g) * av_reg
        aux = {"av_logits": av_logits, "av_reg": av_reg,
               "text_logits": text_logits, "text_reg": text_reg,
               "text_available": text_available, "gate_mean": g.mean()}
        if self.mode == "gate_recon":
            aux["reconstructed_text"] = self.reconstruct(av_h)
        return logits, reg, aux
\end{lstlisting}

\subsubsection*{（3）联合目标与参数更新}
以下节选训练循环中从前向计算到参数更新的部分。
主模型未加入额外的中性类别损失，分类、强度、分支监督和辅助重构
按正文给出的权重联合优化。重构目标由当前训练样本的冻结BERT表示得到。

% Source: src/q2_prediction/experiment_q2_mask_fusion.py :: training loss
\begin{lstlisting}
logits, reg, aux = model(batch, bert)
loss = nn.functional.cross_entropy(logits, batch["class"], weight=weights)
loss += .2 * nn.functional.smooth_l1_loss(reg, batch["strength"])
if args.neutral_loss_weight:
    neutral_logit = logits[:, 1] - torch.logsumexp(logits[:, [0, 2]], dim=-1)
    neutral_target = (batch["class"] == 1).float()
    neutral_count = counts[1]
    neutral_pos_weight = (counts.sum() - neutral_count) / neutral_count
    loss += args.neutral_loss_weight * nn.functional.binary_cross_entropy_with_logits(
        neutral_logit, neutral_target, pos_weight=neutral_pos_weight)
if args.mode in ("gate", "gate_recon"):
    av_loss = nn.functional.cross_entropy(aux["av_logits"], batch["class"], weight=weights)
    av_loss += .2 * nn.functional.smooth_l1_loss(aux["av_reg"], batch["strength"])
    loss += .3 * av_loss
    available = aux["text_available"]
    if available.any():
        text_loss = nn.functional.cross_entropy(aux["text_logits"][available], batch["class"][available], weight=weights)
        text_loss += .2 * nn.functional.smooth_l1_loss(aux["text_reg"][available], batch["strength"][available])
        loss += .2 * text_loss
    if args.mode == "gate_recon":
        # Reconstruction is trained only on clean text samples. The target
        # is the frozen BERT representation, never a validation embedding.
        clean = batch["text_observed"].sum(1) == (
            batch["struct"].sum(1) - 2)
        if clean.any():
            with torch.no_grad():
                h = bert(input_ids=batch["tokens"][clean],
                         attention_mask=batch["struct"][clean].long(),
                         token_type_ids=batch["segment"][clean]).last_hidden_state.float()
                w = batch["text_observed"][clean].unsqueeze(-1).float()
                target = torch.cat([h[:, 0], (h * w).sum(1) / w.sum(1).clamp_min(1)], -1)
            loss += .02 * nn.functional.mse_loss(aux["reconstructed_text"][clean], target)
optimizer.zero_grad(set_to_none=True)
loss.backward()
nn.utils.clip_grad_norm_(model.parameters(), 1.)
optimizer.step()
\end{lstlisting}

\subsubsection*{（4）三种子集成与强度输出}
三个固定模型的分类logits和原始强度分别取平均，再计算类别概率。
类别确定后，强度投影到对应符号区间。

% Source: src/q3_explanation/explain_q3_gate_ensemble.py :: predict
\begin{lstlisting}
def predict(variants, models, bert, chunk=32):
    outputs = []
    with torch.inference_mode():
        for first in range(0, len(variants), chunk):
            group = variants[first:first + chunk]
            batch = {key: torch.cat([row[key] for row in group], dim=0) for key in group[0]}
            preds = [model(batch, bert) for model in models]
            logits = torch.stack([p[0] for p in preds]).mean(0)
            strength = torch.stack([p[1] for p in preds]).mean(0)
            probabilities = torch.softmax(logits, dim=-1)
            for p, z in zip(probabilities.cpu().numpy(), strength.cpu().numpy()):
                outputs.append((p.astype(float), float(z)))
    return outputs
\end{lstlisting}

% Source: src/common/prep.py :: consistent_strength
\begin{lstlisting}
def consistent_strength(classes: np.ndarray, strength: np.ndarray) -> np.ndarray:
    """Project strength onto the official class/sign convention.

    负向取 min(s, -1e-6)、中性取 0、正向取 max(s, 1e-6)，再 clip 到 [-3, 3]。
    """
    return np.where(classes == 1, 0.0,
                    np.where(classes == 0, np.minimum(strength, -1e-6),
                             np.maximum(strength, 1e-6))).clip(-3, 3)
\end{lstlisting}

\subsection{模态贡献与局部证据}
\subsubsection*{（1）模态遮挡}
解释阶段仅修改输入和观测掩码。文本位置替换为掩码词元后重新编码，
语音与视频位置置零；模型权重保持固定。

% Source: src/q3_explanation/explain_q3_gate_ensemble.py :: hide
\begin{lstlisting}
def hide(batch, name, positions):
    key = f"{name}_observed"
    positions = torch.as_tensor(positions, dtype=torch.long, device=batch["tokens"].device)
    observed = batch[key][0, positions]
    positions = positions[observed]
    batch[key][0, positions] = False
    if name == "text":
        batch["tokens"][0, positions] = 103
    else:
        batch[name][0, positions] = 0
\end{lstlisting}

\subsubsection*{（2）八组合贡献分解}
二进制掩码的三位依次表示文本、语音和视频。程序枚举八种输入组合，
固定完整输入的预测类别，计算三模态Shapley值并检查加和关系。

% Source: src/q3_explanation/explain_q3_gate_ensemble.py :: coalition_explanation
\begin{lstlisting}
def coalition_explanation(batch, models, bert):
    variants = []
    masks = list(range(8))
    for mask in masks:
        variant = one(batch, 0)
        for bit, name in enumerate(NAMES):
            if not mask & (1 << bit):
                hide(variant, name, observed_positions(variant, name))
        variants.append(variant)
    predictions = predict(variants, models, bert)
    complete = predictions[7]
    target = int(np.argmax(complete[0]))
    values = np.array([p[target] for p, _ in predictions])
    strength_values = np.array([z for _, z in predictions])
    # 3-player Shapley weights: |S|!(2-|S|)! / 3!.
    weights = (1 / 3, 1 / 6, 1 / 3)
    effects, strength_effects = [], []
    for bit in range(3):
        terms, reg_terms = [], []
        for mask in masks:
            if mask & (1 << bit):
                continue
            weight = weights[mask.bit_count()]
            terms.append(weight * (values[mask | (1 << bit)] - values[mask]))
            reg_terms.append(weight * (strength_values[mask | (1 << bit)] - strength_values[mask]))
        effects.append(float(sum(terms)))
        strength_effects.append(float(sum(reg_terms)))
    if not np.isclose(sum(effects), values[7] - values[0], atol=1e-6):
        raise AssertionError("Shapley additivity failed")
    positive = np.maximum(effects, 0)
    relative = positive / positive.sum() if positive.sum() > 0 else np.zeros(3)
    strength = float(consistent_strength(np.array([target]), np.array([complete[1]]))[0])
    return {"polarity": target, "strength": strength, "full_probability": float(values[7]),
            "empty_probability": float(values[0]), "shapley_probability": effects,
            "shapley_strength_raw": strength_effects, "positive_share": relative.tolist(),
            "main_modality": NAMES[int(np.argmax(effects))],
            "full_probabilities": complete[0].tolist()}
\end{lstlisting}

\subsubsection*{（3）局部窗口敏感性}
逐窗遮挡观测位置，按预测类别概率下降量选取高分、低分和固定随机窗口。

% Source: src/q3_explanation/explain_q3_gate_ensemble.py :: local_explanation
\begin{lstlisting}
def local_explanation(batch, models, bert, target, name):
    positions = observed_positions(batch, name)
    if not len(positions):
        return None
    width = max(1, round(len(positions) * WINDOW_FRACTION))
    windows = [positions[start:start + width] for start in range(len(positions) - width + 1)]
    variants = []
    for window in windows:
        variant = one(batch, 0)
        hide(variant, name, window)
        variants.append(variant)
    base = predict([one(batch, 0)], models, bert)[0][0][target]
    masked = predict(variants, models, bert)
    drops = np.array([base - p[target] for p, _ in masked])
    high = int(np.argmax(drops))
    low = int(np.argmin(drops))
    rng = np.random.default_rng(RANDOM_SEED + int(batch["tokens"][0].sum().item()) + len(name))
    random_idx = int(rng.integers(len(windows)))
    def record(index):
        return {"start_position": int(windows[index][0]),
                "end_position": int(windows[index][-1]),
                "n_observed_positions": int(len(windows[index])),
                "confidence_drop": float(drops[index])}
    return {"high": record(high), "low": record(low), "random": record(random_idx),
            "n_candidate_windows": len(windows), "all_window_drops": drops.tolist()}
\end{lstlisting}
